\documentclass[11pt]{article}
\usepackage{fontspec}
\setmainfont{Latin Modern Roman}
\usepackage{amsmath,amssymb,mathtools,unicode-math}
\setmathfont{Latin Modern Math}
\usepackage{xurl,hyperref}
\hypersetup{colorlinks=true,urlcolor=blue,pdftitle={Problem statement}}
\newcommand{\sref}[2]{\hyperlink{source:#1:#2}{[S#2]}}
\newcommand{\cataloguescope}[1]{\paragraph{Mathematical question.}#1}
\begin{document}
\section{Research Problems in Function Theory — Problem 3.15}\label{statement}

\subsection{Definitions and mathematical statement}

\paragraph{Shared notation.}$\mathbb N=\{1,2,\ldots\}$, and $\mathbb Z,\mathbb Q,\mathbb R,\mathbb C$ denote the integers, rationals, reals and complex numbers. Any other conventions are specified locally.


A subharmonic function is an upper semicontinuous function with values in $[-\infty,\infty)$, not identically $-\infty$ on a connected component, that satisfies the mean-value inequality on every closed ball in its domain. A real harmonic function is a twice continuously differentiable function $u$ satisfying $\Delta u=0$. Let $D\subset\mathbb C$ be a doubly connected domain with boundary curves $\alpha$ and $\beta$, let $z_0,z_1\in D$, and let $A,B\in\mathbb R$. Boundary values mean the ordinary limiting values from $D$. Let $\mathcal U$ consist of subharmonic functions $u$ on $D$ with boundary values $A$ on $\alpha$, $B$ on $\beta$, for which there is a continuous curve $\gamma:[0,1]\to D\cup\alpha$ with $\gamma(0)=z_1$, $\gamma(1)\in\alpha$, $\gamma([0,1))\subset D$, and $u\leq0$ on $\gamma([0,1))$.
\paragraph{Statement recorded in the supplied source.}
Determine $\sup_{u\in\mathcal U}u(z_0)$, whether it is attained, and the maximizing functions, for the given $D,\alpha,\beta,z_0,z_1,A,B$; the curve may depend on $u$. If $\mathcal U$ is empty, the supremum is $-\infty$. \sref{statement}{2}
\subsection{Short English statement}
Maximize a subharmonic value in a doubly connected domain with fixed constants on the two boundaries and a nonpositive path from a second point to the first boundary.
\subsection{Sources}
\begin{description}
\item[\hypertarget{source:statement:1}{S1}] ULAM.AI and the UnsolvedMath Contributors, UnsolvedMath: A Curated Collection of Open Mathematics Problems, catalogue entry. CC BY 4.0. \url{https://huggingface.co/datasets/ulamai/UnsolvedMath}.
\item[\hypertarget{source:statement:2}{S2}] W. K. Hayman and E. F. Lingham, \emph{Research Problems in Function Theory}, arXiv:1809.07200v2 (2018), Problem 3.15. \url{https://arxiv.org/abs/1809.07200}.
\end{description}
\label{end:statement}
\end{document}
